% Este archivo se incorpora desde el documento principal después de
% ejecutar \appendix. No contiene una nueva declaración \appendix.
\chapter{Trilateración acústica por intersección de esferas}
\label{app:trilateracion_esferas}

\section{Alcance y supuestos}

La trilateración por intersección de esferas estima la posición de una
fuente a partir de distancias absolutas medidas desde varios receptores
de posición conocida. Se presenta como una extensión geométrica
exploratoria de los modelos de localización desarrollados en esta tesis.
No corresponde a un método implementado ni validado experimentalmente
en WALO y, por tanto, no se utiliza para producir los resultados
presentados en los capítulos experimentales.

Se considera una fuente en la posición tridimensional
$\mathbf{x}=(x,y,z)$ y $N$ receptores en posiciones conocidas
$\mathbf{r}_i=(x_i,y_i,z_i)$. Si el instante de emisión
$t_{\mathrm{em}}$ es conocido y los retardos instrumentales de cada
canal han sido calibrados, la distancia fuente--receptor puede
estimarse mediante

\begin{equation}
    \rho_i = c(T)\left[
    t_{\mathrm{arr},i}-t_{\mathrm{em}}-t_{\mathrm{off},i}
    \right].
    \label{eq:app_esfera_distancia}
\end{equation}

La ecuación geométrica de cada medición es

\begin{equation}
    \|\mathbf{x}-\mathbf{r}_i\|=\rho_i,
    \label{eq:app_esfera_general}
\end{equation}

o, de forma equivalente,

\begin{equation}
    (x-x_i)^2+(y-y_i)^2+(z-z_i)^2=\rho_i^2.
    \label{eq:app_esfera_expandida}
\end{equation}

Cada medición define una esfera centrada en el receptor. La posición
buscada pertenece a la intersección de todas las superficies
compatibles.

\section{Progresión geométrica según el número de receptores}

\paragraph{Un receptor: una esfera.}
Con un único receptor, la ecuación
\eqref{eq:app_esfera_general} define una esfera de radio $\rho_1$.
La fuente puede encontrarse en cualquier punto de esa superficie, por
lo que la posición tridimensional no queda determinada de manera única.

\paragraph{Dos receptores: una circunferencia.}
Con dos receptores, la posición debe pertenecer simultáneamente a dos
esferas. Si estas se intersectan transversalmente, su intersección es
una circunferencia. El centro de esa circunferencia se encuentra sobre
la línea que une los centros de las esferas. Si las esferas son
tangentes, la intersección se reduce a un punto; si no se intersectan,
las distancias medidas son geométricamente incompatibles en el modelo
ideal.

\paragraph{Tres receptores: dos puntos posibles.}
Una tercera esfera intersecta, en el caso genérico, la circunferencia
definida por las dos primeras. Si existe intersección transversal,
pueden quedar dos puntos posibles, situados a lados opuestos del plano
que contiene los tres centros de los receptores. Esta ambigüedad
refleja que tres centros coplanares no permiten distinguir por sí solos
entre las dos posiciones reflejadas respecto de ese plano. En casos
tangenciales la intersección puede reducirse a un único punto, y en
casos incompatibles puede no existir intersección exacta.

\paragraph{Cuatro receptores: una solución puntual bajo geometría adecuada.}
Una cuarta medición aporta una restricción adicional. Si los cuatro
centros de los receptores no son coplanares, las distancias son
compatibles y la geometría está suficientemente bien condicionada,
la intersección puede determinar una posición única. Cuatro
receptores no garantizan por sí solos unicidad: la disposición
espacial, la incertidumbre de las distancias y la sincronización
determinan si la solución es única y estable.

\section{Formulación algebraica}

Para mostrar cómo se obtiene una solución a partir de las distancias,
se eleva al cuadrado la ecuación~\eqref{eq:app_esfera_general}:

\begin{equation}
    \|\mathbf{x}-\mathbf{r}_i\|^2=\rho_i^2.
\end{equation}

Al restar la ecuación del primer receptor de la ecuación del receptor
$i$, se elimina el término cuadrático $\|\mathbf{x}\|^2$ y se obtiene

\begin{equation}
    2(\mathbf{r}_i-\mathbf{r}_1)^\mathsf{T}\mathbf{x}
    =
    \|\mathbf{r}_i\|^2-\|\mathbf{r}_1\|^2
    +\rho_1^2-\rho_i^2.
    \label{eq:app_esfera_linealizada}
\end{equation}

Estas ecuaciones son lineales en las coordenadas de la fuente. En
tres dimensiones se requieren suficientes restricciones linealmente
independientes para determinar las tres coordenadas. Con cuatro
receptores no coplanares, las tres ecuaciones obtenidas al restar la
primera medición de las otras tres pueden determinar una solución
algebraica, siempre que la matriz geométrica tenga rango completo.
La solución debe verificarse sustituyéndola en las ecuaciones
esféricas originales.

En presencia de ruido, las esferas medidas generalmente no se
intersectan en un punto exacto. En ese caso, puede formularse una
estimación no lineal de mínimos cuadrados:

\begin{equation}
    \hat{\mathbf{x}}=
    \underset{\mathbf{x}}{\operatorname{arg\,min}}
    \sum_{i=1}^{N}w_i
    \left(
    \|\mathbf{x}-\mathbf{r}_i\|-\rho_i
    \right)^2,
    \label{eq:app_esfera_minimos_cuadrados}
\end{equation}

donde $w_i$ representa un peso asociado a la incertidumbre de cada
distancia. La selección de pesos y el análisis de incertidumbre
requieren caracterizar experimentalmente cada canal y la geometría
del arreglo.

\section{Diferencia respecto de TDoA y limitaciones experimentales}

La trilateración por esferas requiere distancias absolutas. Para
obtenerlas a partir de ToA se necesita conocer el instante de emisión,
mantener una referencia temporal común y calibrar los retardos
instrumentales. Si el instante de emisión es desconocido, las
distancias absolutas no pueden obtenerse directamente de los tiempos
de llegada sin introducir una incógnita adicional.

TDoA elimina el instante común de emisión al restar los tiempos
registrados por dos receptores. Esa resta conduce a una diferencia de
distancias constante y, por ello, a una hipérbola en dos dimensiones
o a un hiperboloide en tres dimensiones. PDoA puede representar una
restricción análoga a partir de la diferencia de fase, pero la
periodicidad módulo $2\pi$ introduce ambigüedad en el número de ciclos.
Por tanto, la intersección de esferas y la multilateración hiperbólica
son formulaciones geométricas relacionadas, pero parten de observables
y requisitos de sincronización diferentes.

La viabilidad de este método en WALO requeriría, como mínimo,
incorporar más receptores con posiciones tridimensionales conocidas,
caracterizar los retardos de cada canal, disponer de una referencia
de emisión común y evaluar la incertidumbre de las distancias
estimadas. También sería necesario estudiar la geometría del arreglo,
la influencia de las reflexiones y la estabilidad de la velocidad
del sonido. Estos requisitos exceden la validación experimental
realizada en esta tesis; el método se conserva, por ello, como una
posible línea de extensión y no como una capacidad demostrada del
instrumento.
